\section{Conclusions}
Local chemical environments explain a substantial part of charge-definition disagreement beyond the tested element and source-charge mapping. First-shell information accounts for most of the predictive improvement, with smaller additional gains from two-hop paths. The all-atom and nonmetal result persists when identical source MOFkeys or reduced formulas are blocked, including the records in nonsingleton components.

The chemical interpretation depends on the quantity being summarized. O-attached carbon exhibits a DDEC6--Bader directional shift without an established increase in absolute departure; N-attached carbon exhibits a CM5--Bader magnitude increase despite a near-zero signed contrast. Reporting signed departure, absolute departure, and residual spread separately therefore reveals structure that a single disagreement score would conceal.

The scope is equally specific. Local features worsen the global conditional DDEC6--Bader prediction at metal centers, and a finite cutoff test exposes ambiguous shared-proton contacts even while the listed contrasts remain stable. The released atlas and reproduction materials make both the chemical patterns and these limitations inspectable. They provide a basis for examining charge mappings in particular environments; application-specific validation is needed to connect those patterns to adsorption or other physical observables.
